\section{Open questions}

The five questions below are the residual unknowns after this replication --- 
each is answerable with free-endpoint compute and does not require rerunning
the paper's full study.

\begin{enumerate}
  \item \textbf{Non-Pauli noise (amplitude damping / leakage).}
        \emph{Basis:} RC's argument requires that twirling turns coherent noise
        into stochastic \emph{Pauli} noise, at which point ZNE's linearity
        assumption is well-founded. Amplitude damping twirls to a non-Pauli
        stochastic channel and leakage is not twirled at all by the
        computational Pauli group. Untouched by both paper and this
        reproduction.
        \emph{Next steps:} Add Kraus amplitude damping ($\gamma\in\{10^{-3},
        5\!\times\!10^{-3},10^{-2}\}$) and a 3-level qutrit leakage channel
        ($p_{\text{leak}}\!\sim\!10^{-3}$) per CX; rerun the four-way sweep;
        check whether RC+ZNE still stays flat vs.~$\varepsilon$. Local
        Qiskit-Aer + one free Argo judge.

  \item \textbf{Stacking with virtual distillation (VD, M=2).}
        \emph{Basis:} Our RC+ZNE result plateaus at $\sim\!1.1$--$1.2$\,mHa
        across all $\varepsilon$, apparently limited by the small depolarizing
        residual $+$ finite $N_{\text{rand}}=30$. VD suppresses stochastic
        error at the cost of a duplicate register and a Bell-basis measurement;
        if the residual is genuinely stochastic (as RC should have made it),
        VD should reduce it further.
        \emph{Next steps:} Wrap RC+ZNE in a 2-copy VD executor on 4 qubits,
        estimate $\langle H\otimes I + I\otimes H\rangle/\langle
        \text{SWAP}\rangle$; if the error at $\varepsilon=0.10$ drops
        $\geq\!3\times$ from 1.23\,mHa, stacking is genuinely additive; if not,
        the residual is RC-limited.

  \item \textbf{Transfer to chemically-realistic molecules (BeH$_2$, H$_2$O,
        Fe--S fragment).}
        \emph{Basis:} H$_2$/STO-3G at 2 qubits has trivial entanglement
        structure; deeper problems accumulate coherent error over many more
        CXs. The paper claims LiH also works but LiH was not tested here.
        \emph{Next steps:} BeH$_2$ in 6-qubit parity-tapered mapping, UCC-SD
        ansatz, four-way sweep at $\varepsilon\in\{0.02,0.05,0.10\}$; report
        reduction ratio vs.~CX count.

  \item \textbf{Real-hardware wall-clock overhead vs.~accuracy.}
        \emph{Basis:} $N_{\text{rand}}\!=\!30$ twirls $\times$ 3 ZNE scale
        factors $= 90$ submissions per energy evaluation. On real hardware
        each submission has 50--500\,ms of fixed overhead (compile, queue,
        calibration lookup). A $90\times$ runtime cost that buys $30\times$
        accuracy is only worthwhile if the coherent-noise regime dominates on
        that device that day; on shot-noise-limited or $T_1$-limited devices
        RC+ZNE could be worse per wall-clock second than $90\times$ more shots
        of a raw circuit.
        \emph{Next steps:} Submit raw / RC / ZNE / RC+ZNE H$_2$ circuits to a
        free IBMQ open-plan device at 8192 shots; plot accuracy per wall-clock
        second.

  \item \textbf{Adaptive $(N_{\text{rand}},\text{scale-factor})$ allocation.}
        \emph{Basis:} The static (30, [1,2,3]) budget is inherited from the
        paper. At small $\varepsilon$ RC may be over-invested; at large
        $\varepsilon$ a 4th or 5th scale factor may help. An adaptive scheme
        (allocate more $N_{\text{rand}}$ only if inter-twirl variance is
        large) could match accuracy at half the budget --- or improve it at
        the same budget.
        \emph{Next steps:} Grid $N_{\text{rand}}\in\{5,10,20,30,60\}$ $\times$
        scale-factor sets $\{[1,2], [1,2,3], [1,2,3,5], [1,3,5]\}$ at
        $\varepsilon\in\{0.02,0.10\}$; fit an accuracy-vs-budget Pareto
        frontier; check whether the paper's static point sits on it.
\end{enumerate}
